\documentclass[12pt, a4paper]{article}
\usepackage[utf8]{inputenc}
\usepackage[T2A]{fontenc}
\usepackage{indentfirst, setspace}
\usepackage{tabularx, multirow}
\usepackage[normalem]{ulem}
\usepackage[style=russian]{csquotes}
\usepackage[english,russian]{babel}
\usepackage{hyperref}
\usepackage{ragged2e}
\usepackage{caption}
\usepackage{wrapfig}
\usepackage{amsmath}
\usepackage{tikz}
\makeatletter
\def\@biblabel#1{#1. }
\makeatother
\captionsetup{labelsep=endash}
\usepackage{listings}
\linespread{1.3}
\lstset{
  language=C++, % Язык программирования (можно выбрать другой)
  basicstyle=\ttfamily\small, % Основной стиль шрифта
  keywordstyle=\color{blue}, % Стиль ключевых слов
  breaklines=true,
  commentstyle=\color{green}, % Стиль комментариев
  stringstyle=\color{red}, % Стиль строк
  extendedchars=\true,
  showstringspaces=false,
}

\usepackage[left=20mm, top=15mm, right=15mm, bottom=20mm]{geometry}

\begin{document}
\begin{titlepage}
     \fontsize{12}{12}\selectfont

  {\centering

   \begin{bf}

    \begin{wrapfigure}{l}{19mm}
        
\includegraphics[width=22mm]{emblem.jpeg}
    \end{wrapfigure}


    \noindent Министерство науки и высшего образования Российской Федерации

    \noindent Федеральное государственное бюджетное образовательное учреждение высшего образования

    \noindent \enquote{Московский государственный технический унверситет

     \noindent имени Н.Э. Баумана

     \noindent (национальный исследовательский университет)}

    \noindent (МГТУ им. Н.Э. Баумана)

   \end{bf}
  }

  \vspace{0.4cm}

  {\setstretch{0.1}
   \noindent\rule{\textwidth}{1mm}
   \noindent\rule{\textwidth}{0.5mm}

  }

  \fontsize{14}{21}\selectfont

  \noindent\begin{tabularx}{\textwidth}{l >{\centering\arraybackslash}X}
   ФАКУЛЬТЕТ & \flqq Фундаментальные Науки\frqq \\ \cline{2-2}

   КАФЕДРА & ФН-12 \flqq Математическое моделирование\frqq \\ \cline{2-2}
  \end{tabularx}


  \vspace{1cm}


  \begin{center}
   \begin{bf}

    \fontsize{24}{36}\selectfont
    ОТЧЕТ

    \fontsize{20}{30}\selectfont
    ПО ЛАБОРАТОРНОЙ РАБОТЕ НА ТЕМУ:

    Реализация стека и дека

   \end{bf}
  \end{center}

  \fontsize{14}{21}\selectfont
  \vspace{5cm}


  \noindent\begin{tabularx}{\textwidth}{ X >{\centering}p{4cm} p{1cm} c }
   Студент: & & & Подобедов В. В. \\ \cline{2-2} \cline{4-4}
   & \fontsize{10}{15}\selectfont дата, подпись & & \fontsize{10}{15}\selectfont Ф.И.О. \\
   Преподаватель: & & & Волкова Л. Л.\\ \cline{2-2} \cline{4-4}
   & \fontsize{10}{15}\selectfont дата, подпись & & \fontsize{10}{15}\selectfont Ф.И.О.
  \end{tabularx}

  \vspace{\fill}

  \begin{center}
   \it{Москва}, 2023
  \end{center}

  \thispagestyle{empty}
\end{titlepage}

\renewcommand{\contentsname}{Содержание}
\setcounter{page}{2}
\tableofcontents

\newpage
\begin{center}
   \section*{Введение} 
\addcontentsline{toc}{section}{Введение}
\end{center}

\justifying
\textbf{Цель лабораторной работы}: разработать стек и дек.

Для достижения поставленной цели требуется решить следующие \textbf{задачи}.
\begin{enumerate}
\item Описать две схемы хранения данных в стеке и деке.
\item Разработать алгоритмы, необходимые для хранения данных.
\item Реализовать разработанные алгоритмы.
\item Выполнить тестирование реализации разработанных алгоритмов.
\end{enumerate}

Согласно варианту, требуется разработать следующие две схемы.
\begin{enumerate}
\item Стек. Написать функции добавления элементов, удаления элементов и вывод содержимого в консоль.
\item Дек. разработать функции добавления, удаления данных и их вывод.
\end{enumerate}

\newpage
\begin{center}
   \section{Аналитическая часть} 
\end{center}

\textbf{Стек}\cite{book} --- абстрактный тип данных, представляющий собой список элементов, организованных по принципу LIFO (англ. last in — first out, «последним пришёл — первым вышел»). Стек является очередью с одной стороной входа и выхода элементов. 

\textbf{Дек}\cite{book} --- структура данных типа «список», функционирующая одновременно по двум принцам организации данных: LIFO и FIFO (англ. first in — first out, «первым пришёл — первым вышел»). Определить дек можно как очередь с двумя сторонами, так и стек, имеющий два конца. То есть данный подвид списка характерен двухсторонним доступом: выполнение поэлементной операции, определенной над деком, предполагает возможность выбора одной из его сторон в качестве активной.

\newpage
\begin{center}
    \section{Конструкторская часть}
\end{center}

\subsection{Структура стека} 
Стек состоит из ячеек, которые представлены в виде структуры, содержащей какие-либо данные и указатель на следующий элемент. Для удобного восприятия стек представлен на рисунке ~\ref{img:stack}. 
\begin{figure}[h!]
    \center{\includegraphics[width=10cm, height=10cm]{stack.jpeg}}
    \caption{Схематичное представление стека}
    \label{img:stack}
\end{figure}

Как можно увидеть стек хранит в одном блоке данные и ``*next'' --- ссылку на последующий элемент. Иными словами в ячейке ``*next'' хранится адрес на следующую ячейку.

Над стеком и его элементами можно выполнять следующие операции:
\begin{itemize}
    \item[---] Добавление элемнта в стек.
    \item[---] Удаление элемента из стека.
    \item[---] Проверка стека на пустоту.
    \item[---] Определение количства элементов в стеке.
    \item[---] Вывод всего стека.
\end{itemize}

\subsubsection{Добавление элемента}
Если в стеке еще не было элементов, то создается блок, в котором в ячейку ``данные'' заносится необходимая информация для хранения. Поскольку после этого блока нет больше следующих блоков, то ячейка ``*next'' имеет нулевой указатель. При добавении блока, если уже были созданы блоки, создается новый блок и ячейка ``*next'' в предыдущем блоке начинает ссылаться на созданный блок, а в ячейку ``данные'' добавляется информация.

\subsubsection{Удаление элемента}
При удалении данных, предпоследний блок изменяет ячейку ``*next'' на нулевой указатель, а данные из последнего блока удаляются.

\subsubsection{Проверка стека на пустоту}
Чтобы проверить стек на наличие элементов, необходимо посчитать количество блоков в стеке. Если количество будет больше единицы, то стек не пустой.

\subsubsection{Количество элементов}
Для подсчета количество элементов в стеке нужно создать курсор который будет идти по ячейкам ``*next'' и считать их количество, пока не встретится нулевой указатель.

\subsubsection{Вывод данных}
Аналогично подсчету элементов необходимо создать курсор, который будет идти по ссылкам каждого блока и выводить данные из блока стека.

\subsection{Структура дека} Аналогично стеку содержит блоки в которых хранятся данные и указатель на следующий элемент. Дек схематично изображен на рисунке ~\ref{img:dec}. 
\begin{figure}[h!]
    \center{\includegraphics{dec.png}}
    \caption{Схематичное представление дека}
    \label{img:dec}
\end{figure}

Дек имеет аналогичную структуру, как и стек, в отлчии от того, что добавлять и удалять элементы можно с двух сторон (с конца и с начала). Все остальные функции работают аналогично, как и в стеке.

\subsubsection{Добавление элемента}
При добавлении элемента в дек необходимо написать две функции, которые добавляют элемент либо в начало либо в конец структуры. Если есть необходимость добавить блок в начало то, создается блок и в ячейку ``next'' добавляется ссылка на бывший начальный элемент, а в ячейку ``info'' заносятся необходимые данные. При добавлении блока в конец дека ячейка ``next'' последнего элеменета изменяется на указатель на вновь созданный блок. Данные добавляются в ячейку ``info''.

\subsubsection{Удаление элемнента}
Как и в добавлении элемента в деке можно удалять элемент с любого конца. При удалении элемента из начала дека ячейки ``info'' и ``next'' удаляются методо delete. Если удаляется блок из конца дека, то в предпоследнем блоке ячейку ``next'' необходимо изменить на нулевую ссылку, а последний блок также удалить, воспользовавшись методом delete. 

\newpage
\begin{center}
    \section{Технологическая часть}
\end{center}
Лабораторная работа была реализована на языке программирования C++. Кроме стандартной библиотеки <iostream>, которая предоставляет возможность пользоваься средствами ввода-вывода для стандартной консоли, была подключена библиотека <fstream>. Она включает в себя набор классов, методов и функций, которые предоставляют интерфейс для чтения/записи данных из/в файл. С помощью шаблонных классов ``List'', который описывает саму структуру стека и ``Node'', который описывает блок стека, был реализован стек на ссылках. 

Методы, разработанные в классе, следующие.
\begin{itemize}
    \item[---] Метод ``push\_back'' --- метод, который добавляет элемент в конец стека.
    \item[---] метод ``pop\_back'' --- метод, который удаляет последний элемент из стека.
    \item[---] Метод ``GetSize'' --- метод, который печатает количество элементов в стеке.
    \item[---] Метод ``check'' --- метод, который проверяет стек на пустоту.
    \item[---] Метод ``print'' --- метод, который выводит весь стек в консоль.
\end{itemize}
В программе данные считываются из файла и заносятся в стек. Далее выводится количество элементов, которое хранится в стеке, удаляется последний элемент и проверяется стек на наличие элементов. Если стек не путой, то выводятся данные.

Далее приведен листинг реализации стека (листинг~\ref{lst:label1}).
\begin{lstlisting}[caption={Исходный код программы для реализации стека},label={lst:label1}]
#include <iostream>
#include <fstream>
#include <string>

using namespace std;

template<typename T>
class List {
public:
	List() {
		size = 0;
		head = nullptr;
	}

	~List() {}

	void push_back(T date) {
		if (head == nullptr) {
			head = new Node<T>(date);
		}
		else {
			Node<T> *current = this->head;
			while (current->pNext != nullptr) {
				current = current->pNext;
			}
			current->pNext = new Node<T>(date);
		}
		size++;
	}

	int GetSize() {
		return size;
	}

	void pop_back() {
		Node<T> *previous = this->head;
		for (int i = 0; i < GetSize() - 2; i++){
			previous = previous->pNext;
		}
		Node<T> *toDelete = previous->pNext;
		previous->pNext = toDelete->pNext;
		delete toDelete;
		size--;
	}

	void check() {
		if (size == 0) {
			cout << "Стек пустой" << endl;
		}
		else {
			cout << "В стеке хранятся данные" << endl;
		}
	}

	void print() {
		Node<T> *current = this->head;
		while (current != nullptr) {
			cout << current->date;			
			current = current->pNext;
		}
	}

private:

	template<typename T>
	class Node {
	public:
		Node * pNext;
		T date;

		Node(T date, Node * pNext = nullptr) {
			this->date = date;
			this->pNext = pNext;
		}
	};
	Node<T> *head;
	int size;
};


int main() {
	setlocale(0, "");

	string line;
	ifstream myfile("C:\\Users\\Kononova\\Desktop\\stack_data.txt");

	if (myfile.is_open())
	{
		char c;
		while (myfile.get(c))
		{
			line = line + c;
		}
		myfile.close();
	}
	else
	{
		cout << "Не удалось открыть файл для считывания";
	}

	List<char> lst;

	for (int i = 0; i < line.size(); i++) {
		lst.push_back(line[i]);
	}
	cout << "Количество элементов, включая пробелы: " << lst.GetSize() << endl;
	cout << "Удалили послдний элемент" << endl;
	lst.pop_back();
	cout << "Добавим в конец 5 " << endl;
	lst.push_back('5');
	lst.check();
	if (lst.GetSize() != 0) { 
		lst.print();
	}
	else{
		cout << "стек пустой";
	}

	return 0;
}
\end{lstlisting}

Также, как и в стеке, дек был реализован с помощью шаблонных классов ``List'' и ``Node''. Методы, характерные для дека, следующие.
\begin{itemize}
    \item[---] Метод ``push\_back'' --- метод, который добавляет элемент в конец дек.
    \item[---] Метод ``push\_front'' --- метод, который добавляет элемент в начало дек.
    \item[---] метод ``pop\_back'' --- метод, который удаляет последний элемент из дек.
    \item[---] метод ``pop\_back'' --- метод, который удаляет последний элемент из дек.
    \item[---] Метод ``GetSize'' --- метод, который печатает количество элементов в дек.
    \item[---] Метод ``check'' --- метод, который проверяет дек на пустоту.
    \item[---] Метод ``print'' --- метод, который выводит весь дек в консоль.
\end{itemize}
Далее приведен листинг реализации стека (листинг~\ref{lst:label2}).
\begin{lstlisting}[caption={Исходный код программы для реализации дека},label={lst:label2}]
#include <iostream>
#include <fstream>
#include <string>

using namespace std;

template<typename T>
class List {
public:
	List() {
		size = 0;
		head = nullptr;
	}

	~List() {}

	void push_back(T date) {
		if (head == nullptr) {
			head = new Node<T>(date);
		}
		else {
			Node<T> *current = this->head;
			while (current->pNext != nullptr) {
				current = current->pNext;
			}
			current->pNext = new Node<T>(date);
		}
		size++;
	}

	void push_front(T data) {
		head = new Node<T>(data, head);
		size++;
	}

	int GetSize() {
		return size;
	}

	void pop_back() {
		Node<T> *previous = this->head;
		for (int i = 0; i < GetSize() - 2; i++){
			previous = previous->pNext;
		}
		Node<T> *toDelete = previous->pNext;
		previous->pNext = toDelete->pNext;
		delete toDelete;
		size--;
	}

	void pop_front(){
		Node<T> *temp = head;
		head = head->pNext;
		delete temp;
		size--;
	}

	void check() {
		if (size == 0) {
			cout << "Стек пустой" << endl;
		}
		else {
			cout << "В стеке хранятся данные" << endl;
		}
	}

	void print() {
		Node<T> *current = this->head;
		while (current != nullptr) {
			cout << current->date;			
			current = current->pNext;
		}
	}

private:

	template<typename T>
	class Node {
	public:
		Node * pNext;
		T date;

		Node(T date, Node * pNext = nullptr) {
			this->date = date;
			this->pNext = pNext;
		}
	};
	Node<T> *head;
	int size;
};


int main() {
	setlocale(0, "");

	string line;
	ifstream myfile("C:\\Users\\Kononova\\Desktop\\stack_data.txt");

	if (myfile.is_open())
	{
		char c;
		while (myfile.get(c))
		{
			line = line + c;
		}
		myfile.close();
	}
	else
	{
		cout << "Не удалось открыть файл для считывания";
	}

	List<char> lst;

	for (int i = 0; i < line.size(); i++) {
		lst.push_back(line[i]);
	}
	cout << "Количество элементов, включая пробелы: " << lst.GetSize() << endl;
	cout << "Удалили послдний элемент" << endl;
	lst.pop_back();
	cout << "Удалили первый элемент" << endl;
	lst.pop_front();
	cout << "Добавим в начало r " << endl;
	lst.push_front('r');
	cout << "Добавим в конец e " << endl;
	lst.push_back('e');
	lst.check();
	if (lst.GetSize() != 0) { 
		lst.print();
	}
	else{
		cout << "Дек пустой";
	}

	return 0;
}
\end{lstlisting}

В программе данные считываются из файла и заносятся в дек. Выводится количество элементов, которое хранится в деке, удаляется последний и начальный элемент, добавляется в начало и в конец по элементу и выводятся данные.

\newpage
\begin{center}
    \section{Исследовательская часть}
\end{center}

Далее на рисунках предствалены результаты программы и входной файл, в котором лежат данные.

Результат работы программы для стека приведен на рис. ~\ref{img:grap1}. 
\begin{figure}[h!]
    \center{\includegraphics{stack_1.PNG}}
    \caption{Работоспособность стека}
    \label{img:grap1}
\end{figure}

Результат работы программы для стека приведен на рис. ~\ref{img:grap2}. 
\begin{figure}[h!]
    \center{\includegraphics{stack_2.PNG}}
    \caption{Работоспособность стека}
    \label{img:grap2}
\end{figure}

\newpage
Результат работы программы для дека приведен на рис. ~\ref{img:grap3}. 
\begin{figure}[h!]
    \center{\includegraphics{deck_1.PNG}}
    \caption{Работоспособность дека}
    \label{img:grap3}
\end{figure}

Результат работы программы для дека приведен на рис. ~\ref{img:grap4}. 
\begin{figure}[h!]
    \center{\includegraphics{deck_2.PNG}}
    \caption{Работоспособность дека}
    \label{img:grap4}
\end{figure}

\newpage
\begin{center}
    \section*{Заключение}
    \addcontentsline{toc}{section}{Заключение}
\end{center}
В результате выполнения лабораторной работы были выполнены все задачи.
\begin{enumerate}
\item Описаны две схемы хранения данных в стеке и деке.
\item Разработаны алгоритмы, необходимые для хранения данных.
\item Реализованы разработанные алгоритмы.
\item Выполнено тестирование реализации разработанных алгоритмов.
\end{enumerate}

Согласно варианту, требуется разработаны следующие две схемы.
\begin{enumerate}
\item Стек. Написаны функции добавления элементов, удаления элементов и вывод содержимого в консоль.
\item Дек. разработаны функции добавления, удаления данных и их вывод.
\end{enumerate}

Стек и дек необходимы для хранения информации. Дек, в отличии от стека является более гибким хранилищем информации, так как он располагает большим набором функций. 

\newpage
\begin{center}
\begin{thebibliography}{}
\addcontentsline{toc}{section}{Список используемых источников}
\bibitem{book} Е.А. Конова , Г.А. Поллак, А.М. Ткачев. Структуры данных. Программирование на языке С и С++. Учебное пособие, под ред. Поллак Г.А., Челябинск, Издательство ЮУрГУ, 2004 
\end{thebibliography}
\end{center}
\end{document}
